\usepackage{fontspec}
\usepackage{polyglossia}
\setdefaultlanguage{russian}
\setotherlanguage{english}
% Предпочтительный шрифт по требованиям: Times New Roman.
% Для контейнерной сборки с кириллицей оставлен безопасный fallback.
\IfFontExistsTF{Times New Roman}{
  \setmainfont{Times New Roman}
  \newfontfamily\cyrillicfont{Times New Roman}
}{
  \IfFontExistsTF{FreeSerif}{
    \setmainfont{FreeSerif}
    \newfontfamily\cyrillicfont{FreeSerif}
  }{
    \IfFontExistsTF{Linux Libertine O}{
      \setmainfont{Linux Libertine O}
      \newfontfamily\cyrillicfont{Linux Libertine O}
    }{
      \setmainfont{TeX Gyre Termes}
      \newfontfamily\cyrillicfont{TeX Gyre Termes}
    }
  }
}
\IfFontExistsTF{Courier New}{
  \newfontfamily\cyrillicfonttt{Courier New}
}{
  \IfFontExistsTF{TeX Gyre Cursor}{
    \newfontfamily\cyrillicfonttt{TeX Gyre Cursor}
  }{
    \newfontfamily\cyrillicfonttt{Latin Modern Mono}
  }
}

\usepackage[a4paper,left=30mm,right=15mm,top=20mm,bottom=20mm]{geometry}
\usepackage{setspace}
\onehalfspacing
\setlength{\parindent}{1.25cm}
\setlength{\parskip}{0pt}
\usepackage{indentfirst}

\usepackage{amsmath}
\usepackage{amssymb}
\usepackage{booktabs}
\usepackage{longtable}
\usepackage{array}
\usepackage{tabularx}
\usepackage{multirow}
\usepackage{enumitem}
\setlist[itemize]{
  label=-,
  leftmargin=1.25cm,
  listparindent=1.25cm,
  labelsep=0.5em,
  itemsep=0pt,
  topsep=6pt,
  parsep=0pt,
  partopsep=0pt
}
\setlist[enumerate]{
  leftmargin=1.25cm,
  listparindent=1.25cm,
  labelsep=0.5em,
  itemsep=0pt,
  topsep=6pt,
  parsep=0pt,
  partopsep=0pt
}
\setlist[enumerate,1]{label=\arabic*),ref=\arabic*)}
\setlist[enumerate,2]{label=\alph*),ref=\alph*)}
\setlist[enumerate,3]{label=\roman*),ref=\roman*)}

\usepackage{graphicx}
\graphicspath{{figures/}}
\usepackage{float}
\usepackage{tikz}
\usetikzlibrary{arrows.meta,positioning,shapes.geometric,shapes.symbols,fit,backgrounds,calc,shadows.blur}
\usepackage{xcolor}
\definecolor{cBg}{RGB}{246,247,250}
\definecolor{cBox}{RGB}{233,239,247}
\definecolor{cBoxAlt}{RGB}{241,237,231}
\definecolor{cBoxAcc}{RGB}{226,238,233}
\definecolor{cBorder}{RGB}{120,134,154}
\definecolor{cArrow}{RGB}{60,75,100}
\definecolor{cText}{RGB}{30,33,40}
\definecolor{cCodeBg}{RGB}{248,248,244}
\definecolor{cCodeKw}{RGB}{0,90,160}
\definecolor{cCodeStr}{RGB}{120,40,40}
\definecolor{cCodeCmt}{RGB}{90,110,90}
\definecolor{cCodeNum}{RGB}{120,120,120}
\usepackage{caption}
\DeclareCaptionLabelSeparator{emdash}{ - }
\captionsetup[figure]{name=Рисунок,labelsep=emdash,justification=centering}
\captionsetup[table]{name=Таблица,labelsep=emdash,justification=raggedright,singlelinecheck=false}
\captionsetup{skip=6pt}
\setlength{\abovecaptionskip}{6pt}
\setlength{\belowcaptionskip}{0pt}

\usepackage{chngcntr}
\counterwithin{figure}{chapter}
\counterwithin{table}{chapter}
\renewcommand{\theequation}{\arabic{equation}}

\usepackage[hidelinks]{hyperref}
\usepackage{xurl}
\usepackage{csquotes}
\usepackage[
  backend=biber,
  style=gost-numeric,
  sorting=none,
  language=auto,
  autolang=other,
  bibencoding=utf8
]{biblatex}
\DeclareLanguageMapping{russian}{russian-gost}
\urlstyle{same}

\usepackage{fancyhdr}
\pagestyle{fancy}
\fancyhf{}
\fancyfoot[C]{\thepage}
\renewcommand{\headrulewidth}{0pt}
\renewcommand{\footrulewidth}{0pt}

\setcounter{secnumdepth}{3}
\setcounter{tocdepth}{2}

\usepackage[explicit]{titlesec}
\newcommand{\headinghyphenpenalties}{%
  \hyphenpenalty=10000%
  \exhyphenpenalty=10000%
  \relax%
}
\titleformat{\chapter}[block]
  {\bfseries\centering\headinghyphenpenalties}
  {\thechapter}
  {1em}
  {\MakeUppercase{#1}}
\titleformat{name=\chapter,numberless}[block]
  {\bfseries\centering\headinghyphenpenalties}
  {}
  {0pt}
  {\MakeUppercase{#1}}
\titlespacing*{\chapter}{0pt}{0pt}{20pt}
\titleformat{\section}[block]
  {\bfseries\headinghyphenpenalties}
  {\thesection}
  {1em}
  {#1}
\titlespacing*{\section}{0pt}{0.7\baselineskip}{0.5\baselineskip}
\titleformat{\subsection}[block]
  {\bfseries\headinghyphenpenalties}
  {\thesubsection}
  {1em}
  {#1}
\titlespacing*{\subsection}{0pt}{0.7\baselineskip}{0.5\baselineskip}
\titleformat{\subsubsection}[block]
  {\bfseries\headinghyphenpenalties}
  {\thesubsubsection}
  {1em}
  {#1}
\titlespacing*{\subsubsection}{0pt}{0.7\baselineskip}{0.5\baselineskip}
% Фиксируем название раздела table of contents строго как «Содержание».
% Для polyglossia задаем имя через русские captions и дополнительно в BeginDocument.
\gappto\captionsrussian{\renewcommand{\contentsname}{Содержание}}
\AtBeginDocument{\renewcommand{\contentsname}{Содержание}}

% ГОСТ: формулы с визуально отделенными интервалами, единые интервалы float-элементов.
\setlength{\abovedisplayskip}{\baselineskip}
\setlength{\belowdisplayskip}{\baselineskip}
\setlength{\abovedisplayshortskip}{\baselineskip}
\setlength{\belowdisplayshortskip}{\baselineskip}
\setlength{\textfloatsep}{\baselineskip}
\setlength{\floatsep}{0.8\baselineskip}
\setlength{\intextsep}{0.8\baselineskip}

% Уменьшает переполнения строк без изменения полей и кегля.
\emergencystretch=2em
\tolerance=1600

\usepackage{listings}
\lstdefinestyle{codestyle}{
  basicstyle=\ttfamily\scriptsize,
  backgroundcolor=\color{cCodeBg},
  keywordstyle=\color{cCodeKw}\bfseries,
  stringstyle=\color{cCodeStr},
  commentstyle=\color{cCodeCmt}\itshape,
  numberstyle=\tiny\color{cCodeNum},
  numbers=left,
  numbersep=8pt,
  showstringspaces=false,
  breaklines=true,
  breakatwhitespace=false,
  postbreak=\mbox{\textcolor{cCodeNum}{$\hookrightarrow$}\space},
  columns=fullflexible,
  keepspaces=true,
  frame=single,
  rulecolor=\color{cBorder},
  framesep=4pt,
  xleftmargin=14pt,
  xrightmargin=2pt,
  aboveskip=8pt,
  belowskip=8pt,
  tabsize=2,
  extendedchars=true,
  inputencoding=utf8
}
\lstset{style=codestyle}
\AtBeginDocument{\counterwithin{lstlisting}{chapter}}
\renewcommand{\lstlistingname}{Листинг}

% Оформление пояснений под формулами по ГОСТ 6.8.2.
\newcommand{\where}{\par\noindent\hspace*{\parindent}где\ }
\newcommand{\whereitem}[2]{%
  \ifhmode\else\noindent\hspace*{\parindent}\fi
  #1 \textemdash\ #2;\par
}
\newcommand{\whereitemlast}[2]{%
  \ifhmode\else\noindent\hspace*{\parindent}\fi
  #1 \textemdash\ #2.\par
}
